\documentclass[aspectratio=169,10pt,english]{beamer}
\input{../../../_preambulo_beamer.tex}
\input{../../../_comandos_beamer.tex}
\selectlanguage{english}
\title{L07: Counting Techniques II}
\subtitle{Class route -- 40 minutes}
\author{}
\date{}
\begin{document}

\begin{frame}{Repeated categories and allocations}
\framesubtitle{Class route -- 40 minutes}
\textbf{Question:} What changes when categories repeat, and why can a correct count still give a wrong probability?
\begin{block}{By the end, you can}
Count distinct repeated-label sequences, interpret fixed-size multinomial allocations, count occupancy profiles, and explain why profiles need not be equally likely.
\end{block}
\smallskip
All requests, queues, service centers, and routing rules are synthetic teaching assumptions.
\end{frame}

\begin{frame}{A 40-minute route}
\small
\begin{tabular}{p{.43\linewidth}rp{.38\linewidth}}
\toprule Topic & Min & Observable check \\ \midrule
Define the outcome & 4 & Separate requests and labels.\\
Repeated permutations & 8 & Derive and enumerate 560.\\
Multinomial allocation & 9 & Verify 210--210--140.\\
Stars and bars; probability & 15 & Count 165; reject $1/165$.\\
Integration & 4 & Explain why counts need weights.\\ \bottomrule
\end{tabular}
\par\medskip
Run the prepared cells in \texttt{lesson\_07\_compact.ipynb} alongside these slides. Code takes seconds; the estimate includes derivation and checks. Additional figures and practice remain in the extended materials.
\end{frame}

\begin{frame}{Define what is distinct}
\small
Eight \textbf{labeled requests} $R01,\ldots,R08$ receive queue labels with fixed sizes: three Priority, three Standard, and two Audit.
\par\medskip
In a sequence of \emph{queue labels}, swapping two Priority entries changes nothing. In an allocation of \emph{requests}, $R01$ and $R02$ are different items. We will count both views with the same coefficient but interpret them differently.
\par\medskip
\textbf{Check:} Does exchanging two identical queue labels create a new sequence? No. Does exchanging two labeled requests between queues change the allocation? Yes.
\end{frame}

\begin{frame}{Repeated permutations remove duplicates}
\begin{columns}[T]
\begin{column}{.49\textwidth}
\small
Treating all eight labels as distinct gives $8!=40{,}320$, but each unique sequence appears $3!3!2!=72$ times.
\[
\frac{8!}{3!3!2!}=560.
\]
The notebook enumerates the unique sequences and confirms 560. The denominator removes exchanges of equal labels, not a probability weight.
\par\smallskip
\textbf{Check:} Why is $8!$ too large?
\end{column}
\begin{column}{.49\textwidth}
\centering
\includegraphics[width=\linewidth]{../figures/allocations_01_repeated_permutations.png}
\end{column}
\end{columns}
\end{frame}

\begin{frame}[fragile]{Code check: formula and enumeration}
\small
The two calculations answer the same question about distinct label sequences:
\begin{lstlisting}
from itertools import permutations
from math import factorial
labels = (["Priority"] * 3
          + ["Standard"] * 3
          + ["Audit"] * 2)
formula = factorial(8) // (factorial(3)
                           * factorial(3)
                           * factorial(2))
observed = len(set(permutations(labels)))
assert formula == observed == 560
\end{lstlisting}
Enumeration verifies the model for this small case; the formula explains the general rule.
\end{frame}

\begin{frame}{The same coefficient allocates requests}
\begin{columns}[T]
\begin{column}{.48\textwidth}
\small
Choose three of eight labeled requests for Priority, three of the remaining five for Standard, and the last two for Audit:
\[
\binom83\binom53\binom22=56\times10\times1=560.
\]
Across all 560 allocations, $R01$ is in Priority 210 times, Standard 210 times, and Audit 140 times. These are $3/8$, $3/8$, and $2/8$ of the allocations.
\par\smallskip
\textbf{Check:} The coefficient is the same as before; what changed? The interpretation of one outcome.
\end{column}
\begin{column}{.50\textwidth}
\centering
\includegraphics[width=\linewidth]{../figures/allocations_02_multinomial_categories.png}
\end{column}
\end{columns}
\end{frame}

\begin{frame}{Stars and bars counts occupancy profiles}
\small
Now route eight requests to \textbf{four labeled centers}. A profile $(n_1,n_2,n_3,n_4)$ records how many each center receives, with $n_i\geq0$ and $\sum_i n_i=8$.
\[
\#\{(n_1,n_2,n_3,n_4)\}=\binom{8+4-1}{4-1}=\binom{11}{3}=165.
\]
The requests are labeled, but this calculation counts only the possible \emph{count vectors}. For example, $(8,0,0,0)$ and $(0,8,0,0)$ differ because centers are labeled.
\par\medskip
\textbf{Check:} Does 165 alone tell us the probability of a particular profile? No; it counts possibilities without assigning weights.
\end{frame}

\begin{frame}{Random routing does not make profiles equal}
\small
Suppose every labeled request independently chooses one of four centers with probability $1/4$. For profile $\boldsymbol n=(n_1,n_2,n_3,n_4)$,
\[
P(\boldsymbol n)=\frac{8!}{n_1!n_2!n_3!n_4!}\left(\frac14\right)^8.
\]
The factorial factor counts how many labeled-request assignments produce that profile.
\begin{center}
\begin{tabular}{lrr}
\toprule Profile & Exact $P$ & Incorrect $1/165$ \\ \midrule
$(5,1,1,1)$ & 0.00512695 & 0.00606061\\
$(3,2,2,1)$ & 0.02563477 & 0.00606061\\ \bottomrule
\end{tabular}
\end{center}
\end{frame}

\begin{frame}{Compare calculation and simulation}
\begin{columns}[T]
\begin{column}{.43\textwidth}
\small
The notebook enumerates 165 profiles, computes their selected exact probabilities, and uses 500,000 seed-42 simulations.
\par\medskip
For $(5,1,1,1)$, the simulation gives 0.005288; for $(3,2,2,1)$ it gives 0.025894. These approximate different exact values, not a common $1/165$.
\par\smallskip
\textbf{Check:} Which count is appropriate for a probability numerator? The number of labeled assignments producing the profile, given the routing model.
\end{column}
\begin{column}{.55\textwidth}
\centering
\includegraphics[width=\linewidth]{../figures/allocations_03_star_bars_probability_limit.png}
\end{column}
\end{columns}
\end{frame}

\begin{frame}{Count the right objects; state the probability model}
\small
\begin{enumerate}
\item Repeated labels: divide out exchanges that do not change the sequence.
\item Fixed-size queues: allocate labeled requests with a multinomial count.
\item Occupancy profiles: stars and bars counts possible vectors.
\item Probabilities: specify how requests are routed and how many assignments lead to each vector.
\end{enumerate}
\par\medskip
\textbf{Final check:} Why are there 165 profiles but $(3,2,2,1)$ is much more likely than $(5,1,1,1)$? More labeled-request assignments produce the former under independent uniform routing.
\par\smallskip
\textbf{Next:} L08 develops probability rules for events rather than treating all counted categories as equally likely.
\end{frame}
\end{document}
